\documentclass[10pt,UTF8]{article}


\usepackage{geometry}

\usepackage[scheme=plain]{ctex}

\geometry{a4paper,scale=0.9}

% \usepackage{times}  % 设置正文字体为 Adobe Times Roman
% \usepackage{mathptmx} % 只加载数学字体
% \renewcommand{\rmdefault}{cmr} % 恢复正文字体为 Computer Modern Roman

\usepackage{bm} % 数学粗体（避免斜体被取消）

% \usepackage{makecell}
% \usepackage{multirow} % 表格多行合并

\usepackage{booktabs} % 表格美化

% \usepackage{graphicx} %插入图片的宏包
% \usepackage{subfigure}
% \usepackage{float} %设置图片浮动位置的宏包

\usepackage{amsmath}
\usepackage{esint} % 提供更丰富的积分符号
\usepackage{amssymb}
\usepackage{mathtools}

\usepackage{physics}

% \usepackage{siunitx}
% % 关闭警告
% \ExplSyntaxOn
% \msg_redirect_name:nnn { siunitx } { physics-pkg } { none }
% \ExplSyntaxOff

% \usepackage{bussproofs} % 自然演绎推理树

% 交叉引用
% 关闭边框，设置颜色
\usepackage[colorlinks=true, linkcolor=darkgray, citecolor=teal,
urlcolor=blue]{hyperref}
% 使 \cref 和 \Cref 可用
\usepackage{cleveref}

\usepackage{bookmark}  % 生成 PDF 大纲

\allowdisplaybreaks[3]  % 允许换页


%%%%%%%%%%%%%%%%%%%%%%%% tcolorbox %%%%%%%%%%%%%%%%%%%%%%%%%%
\usepackage[most]{tcolorbox}

\newenvironment{tipBox}
{
  \begin{tcolorbox}[colback=green!5!white,colframe=green!75!black,parbox=false,before
  upper=\par,before lower=\par]}
  {
\end{tcolorbox}}

\newenvironment{conceptBox}
{
\begin{tcolorbox}[colback=blue!5!white,colframe=blue!75!black]}
  {
\end{tcolorbox}}

\newenvironment{theoremBoxNoIndent}
{
\begin{tcolorbox}[colback=yellow!5!white,colframe=yellow!75!black]}
  {
\end{tcolorbox}}

\newenvironment{definitionBox}
{
  \begin{tcolorbox}[colback=blue!5!white,colframe=blue!75!black,parbox=false,before
  upper=\par,before lower=\par]}
  {
\end{tcolorbox}}

\newenvironment{theoremBox}
{
  \begin{tcolorbox}[colback=yellow!5!white,colframe=yellow!75!black,parbox=false,before
  upper=\par,before lower=\par]}
  {
\end{tcolorbox}}
%%%%%%%%%%%%%%%%%%%%%%%%%%%%%%%%%%%%%%%%%%%%%%%%%%%%%%%%%%%%

%%%%%%%%%%%%%%%%%%%%% amsthm %%%%%%%%%%%%%%%%
\usepackage{amsthm}

\theoremstyle{definition}
\newtheorem{definition inner}{Definition}[section]
\newenvironment{definition}[1][]
{
  \begin{definitionBox}
  \begin{definition inner}[#1]\hfill\par}
    {
    \end{definition inner}
\end{definitionBox}}

\theoremstyle{definition}
\newtheorem{theorem inner}{Theorem}[section]
\newenvironment{theorem}[1][]
{
  \begin{theoremBox}
  \begin{theorem inner}[#1]\hfill\par}
    {
    \end{theorem inner}
\end{theoremBox}}

\theoremstyle{definition}
\newtheorem*{proof inner}{\textit{Proof}}
\renewenvironment{proof}[1][]
{
\begin{proof inner}[#1]\hfill\par}
  {\qed\vphantom{.}
\end{proof inner}}

\theoremstyle{plain}
\newtheorem{tip inner}{Tip}[section]
\newenvironment{tip}[1][]
{
  \begin{tipBox}
  \begin{tip inner}[#1]\hfill\par}
    {
    \end{tip inner}
\end{tipBox}}

\theoremstyle{remark}
\newtheorem*{remark}{\textit{Remark}}
%%%%%%%%%%%%%%%%%%%%%%%%%%%%%%%%%%%%%%%%%%%%%%


\newcommand{\iu}{\mathrm{i}}  % 虚数单位 i


\begin{document}

\part*{Second Order ODEs}

The \textbf{general form} of a second order ODE is
\[
  \dv[2]{y}{t} = f(t,y,\dv{y}{t})
\]
for some given function \(f\).

\section{Linear Homogeneous Equations}

\begin{definition}[Linear Equation]
  A second order ODE is called \textbf{linear}
  if it can be written in the form
  \[
    a(t)\cdot y'' + b(t)\cdot y' + c(t)\cdot y = g(t)
  \]
\end{definition}

If \(a(t) \neq 0\), we can divide both sides by \(a(t)\) to get
\( y'' + p(t)\cdot y' + q(t)\cdot y = r(t) \).

\begin{definition}[Homogeneous Equation]
  A linear equation is called \textbf{homogeneous}
  if it can be written in the form
  \[
    a(t)\cdot y'' + b(t)\cdot y' + c(t)\cdot y = 0
  \]
  Otherwise, it is called \textbf{non-homogeneous}.
\end{definition}

\begin{theorem}[Existence and Uniqueness Theorem]
  {}\label{thm: Existence and Uniqueness Theorem}
  Consider the IVP\@:
  \begin{equation*}
    y'' + p(t)\cdot y' + q(t)\cdot y = r(t),
    \qquad y(t_0) = y_0, \quad y'(t_0) = y_1
  \end{equation*}
  if there exists an open interval \(I \subseteq \mathbb{R}\)
  such that \(t_0 \in I\) and \(p\), \(q\) and \(r\)
  are \textbf{continuous} on \(I\),
  then there is a \textbf{exact one} solution \(y = \phi(t)\),
  and it exists \textbf{throughout} \(I\).
\end{theorem}

\begin{theorem}[Principle of Superposition]
  {}\label{thm: Principle of Superposition}
  If \(y_1\) and \(y_2\) are two solutions of the \textbf{homogeneous equation}
  \[
    a(t)\cdot y'' + b(t)\cdot y' + c(t)\cdot y = 0
  \]
  then any linear combination \(c_1 y_1(t) + c_2 y_2(t)\)
  is also a solution of the ODE\@.
\end{theorem}

\begin{definition}[Wronskian]
  {}\label{def: Wronskian}
  The \textbf{Wronskian}, or \textbf{Wronskian determinant},
  of two functions \(y_1\) and \(y_2\) is defined as
  \[
    W(y_1,y_2)(t) :=
    \begin{vmatrix}
      y_1(t) & y_2(t) \\
      y_1'(t) & y_2'(t)
    \end{vmatrix} = y_1(t)\cdot y_2'(t) - y_2(t)\cdot y_1'(t)
  \]
\end{definition}

\begin{theorem}[Abel's Theorem]
  {}\label{thm: Abel's Theorem}
  Suppose \(p(t)\) and \(q(t)\) are continuous on an open interval \(I\),
  and \(y_1\) and \(y_2\) are two non-zero solutions of the ODE
  \begin{equation}
    y'' + \textcolor{blue}{p(t)}\cdot y' + q(t)\cdot y = 0
    \label{eq: local.06 ODE}
  \end{equation}
  Then, the Wronskian of \(y_1\) and \(y_2\) is given by
  \begin{equation*}
    W(y_1,y_2)(t) = \textcolor{red}{c} \cdot
    \exp(-\int \textcolor{blue}{p(t)} \dd{t})
    \qquad\text{ where } \textcolor{red}{c} \in \mathbb{R}
    \text{ is a constant depending on } y_1, y_2
  \end{equation*}
\end{theorem}

\begin{remark}
  If \(y(t),p(t),q(t) \in \mathbb{C}\),
  then \(\textcolor{red}{c} \in \mathbb{C}\).
\end{remark}

\begin{theoremBox}\noindent
  \textbf{Corollary}
  {}\label{thm: corollary of Abel's Theorem}
  \qquad\qquad\qquad
  \begin{math}
    \exists t_0 \in I \text{ s.t. } W(y_1,y_2)(t_0) \neq 0
    \quad\implies\quad
    \forall t \in I, \; W(y_1,y_2)(t) \neq 0
  \end{math}
\end{theoremBox}

\begin{proof}[Abel's Theorem]\noindent
  \begin{align*}
    \begin{cases}
      y_2 y_1'' + y_2 p y_1' + y_2 q y_1 = 0 \\
      y_1 y_2'' + y_1 p y_2' + y_1 q y_2 = 0
    \end{cases}
    &\implies
    (y_2 y_1'' - y_1 y_2'') + p \cdot (y_2 y_1' - y_1 y_2') = 0 \\
    &\implies
    W' + p \cdot W = 0 \\
    &\implies
    \abs{W(t)} = \exp(-\int p(t) \dd{t})
    \qor W(t) = 0
  \end{align*}
\end{proof}

\begin{theorem}[Linear Dependence and Wronskian]
  {}\label{thm: linear dependence and Wronskian}
  Suppose \(y_1\) and \(y_2\) are two solutions of an ODE
  in the form of \cref{eq: local.06 ODE}. Then,
  \begin{equation*}
    y_1(t) \qand y_2(t) \qq{are \textbf{linearly independent} on}
    t \in I
    \qquad\iff\qquad
    \forall t \in I, \;
    W(y_1,y_2)(t) \neq 0
  \end{equation*}
\end{theorem}

\begin{proof}\noindent
  \begin{align*}
    \text{linear independence}
    \quad&\iff\quad
    \forall c_1, c_2 \in \mathbb{C},\;
    \Big(c_1 y_1(t) + c_2 y_2(t) = 0\Big)
    \; \to \; \Big(c_1 = c_2 = 0\Big)
    \tag*{[Proposition 1]} \\
    &\iff\quad
    \forall c_1, c_2 \in \mathbb{C},\;
    \begin{bmatrix}
      y_1(t) & y_2(t) \\
      y_1'(t) & y_2'(t)
    \end{bmatrix}
    \begin{bmatrix}
      c_1 \\ c_2
    \end{bmatrix} =
    \begin{bmatrix}
      0 \\ 0
    \end{bmatrix} \; \to \;
    \begin{bmatrix}
      c_1 \\ c_2
    \end{bmatrix} =
    \begin{bmatrix}
      0 \\ 0
    \end{bmatrix}
    \tag*{[Proposition 2]} \\
    &\iff\quad
    W(y_1,y_2)(t) \not\equiv 0 \\
    &\iff\quad
    \forall t \in I, \; W(y_1,y_2)(t) \neq 0
    \qq{by the \hyperref[thm: corollary of Abel's Theorem]{
    corollary of Abel's Theorem}}
  \end{align*}
  \big([Proposition 2] \(\implies \) [Proposition 1]\big) is obtained by
  combining [Proposition 2] with the following fact:
  \begin{equation*}
    \forall c_1, c_2 \in \mathbb{C},\;
    \Big(c_1 y_1(t) + c_2 y_2(t) = 0\Big)
    \; \to \;
    \Big(c_1 y_1'(t) + c_2 y_2'(t) = 0\Big)
  \end{equation*}
\end{proof}

\begin{theorem}
  {}\label{thm: Wronskian and IVP with homogeneous ODE}
  Suppose \(p(t)\) and \(q(t)\) are continuous on an open interval \(I\),
  and \(y_1\) and \(y_2\) are two solutions of the ODE
  \begin{equation}
    y'' + p(t)\cdot y' + q(t)\cdot y = 0
    \qfor t \in I
    \label{eq: local.01 ODE of y}
  \end{equation}
  Let \(\phi(t)\) be the \textbf{unique} solution
  (by Theorem~\ref{thm: Existence and Uniqueness Theorem})
  of the IVP\@:
  \begin{equation}
    y'' + p(t)\cdot y' + q(t)\cdot y = 0
    \qquad y(t_0) = m, \quad y'(t_0) = n
    \qfor t \in I \qand t_0 \in I
    \label{eq: local.02 IVP of y}
  \end{equation}
  Then, for \(t \in I\),
  \begin{equation}
    \exists !\, c_1, c_2 \in \mathbb{R} \text{ s.t. }
    \phi(t) = c_1 y_1(t) + c_2 y_2(t)
    \qquad\iff\qquad
    W(y_1,y_2)(t_0) \neq 0
    \label{eq: local.05 the proposition}
  \end{equation}
\end{theorem}

\begin{remark}
  In the assumption, \(y_1\) and \(y_2\) are not required to be
  linearly independent.
  Thus \(c_1,c_2\) are required to be unique in
  the left hand side of
  \cref{eq: local.05 the proposition}.
\end{remark}

\begin{proof}\noindent
  \subparagraph*{Sufficiency.}
  Since \(\phi(t_0) = m\) and \(\phi'(t_0) = n\), we have
  \begin{align*}
    \exists !\, c_1, c_2 \in \mathbb{R} \text{ s.t. }
    \phi(t) = c_1 y_1(t) + c_2 y_2(t)
    \quad&\implies\quad
    \exists !\, c_1, c_2 \in \mathbb{R}\;
    \text{ s.t.}
    \begin{bmatrix}
      y_1(t_0) & y_2(t_0) \\
      y_1'(t_0) & y_2'(t_0)
    \end{bmatrix}
    \begin{bmatrix} c_1 \\ c_2
    \end{bmatrix} =
    \begin{bmatrix} m \\ n
    \end{bmatrix} \\
    &\implies\quad
    \exists b_1, b_2 \in \mathbb{R},\;
    \exists !\, c_1, c_2 \in \mathbb{R}\;
    \text{ s.t.}
    \begin{bmatrix}
      y_1(t_0) & y_2(t_0) \\
      y_1'(t_0) & y_2'(t_0)
    \end{bmatrix}
    \begin{bmatrix} c_1 \\ c_2
    \end{bmatrix} =
    \begin{bmatrix} b_1 \\ b_2
    \end{bmatrix} \\
    &\iff\quad
    W(y_1,y_2)(t_0) \neq 0
  \end{align*}
  \subparagraph*{Necessity.}
  \begin{equation*}
    W(y_1,y_2)(t_0) \neq 0
    \quad\implies\quad
    \exists !\, c_1, c_2 \in \mathbb{R}\;
    \text{ s.t.}
    \begin{bmatrix}
      y_1(t_0) & y_2(t_0) \\
      y_1'(t_0) & y_2'(t_0)
    \end{bmatrix}
    \begin{bmatrix}
      c_1 \\ c_2
    \end{bmatrix} =
    \begin{bmatrix}
      m \\ n
    \end{bmatrix}
  \end{equation*}
  Let \(\varphi(t) := c_1 y_1(t) + c_2 y_2(t), t\in I\).
  Then we have \(\varphi(t_0) = m\) and \(\varphi'(t_0) = n\).
  \begin{gather*}
    \text{Theorem~\ref{thm: Principle of Superposition}
    (\nameref{thm: Principle of Superposition})}
    \quad\implies\quad
    \varphi(t) \text{ is a solution of the ODE \cref{eq: local.01 ODE of y}.} \\
    \text{Theorem~\ref{thm: Existence and Uniqueness Theorem}
    (\nameref{thm: Existence and Uniqueness Theorem})}
    \quad\implies\quad
    \varphi(t) \text{ is the \textbf{only} solution of the IVP in
    \cref{eq: local.02 IVP of y}.}
  \end{gather*}
  Thus, \(\varphi(t) = \phi(t)\) and we have
  \begin{equation*}
    W(y_1,y_2)(t_0) \neq 0
    \quad\implies\quad
    \exists !\, c_1, c_2 \in \mathbb{R} \text{ s.t. }
    \phi(t) = \varphi(t) = c_1 y_1(t) + c_2 y_2(t)
  \end{equation*}
\end{proof}

\begin{theorem}
  {}\label{thm: Wronskian and general solution of a homogeneous ODE}
  Suppose \(p(t)\) and \(q(t)\) are continuous on an open interval \(I\),
  and \(y_1\) and \(y_2\) are two solutions of ODE
  \begin{equation*}
    y'' + p(t)\cdot y' + q(t)\cdot y = 0
  \end{equation*}
  for \(t \in I\).
  Then,
  \begin{equation*}
    \text{The \textbf{general solution} is given by }
    \phi(t) = c_1 y_1(t) + c_2 y_2(t).
    \quad\iff\quad
    \exists t_0 \in I \text{ s.t. } W(y_1,y_2)(t_0) \neq 0.
  \end{equation*}

  \(y_1, y_2\) are said to form a
  \textbf{fundamental set of solutions (FSS)} of the ODE,
  if their Wronskian is non-zero.
\end{theorem}

\begin{remark}
  Theorem~\ref{thm: Wronskian and general solution of a homogeneous ODE}
  is used to determine whether a given pair of solutions
  \(y_1, y_2\) form a FSS,
  rather than verifying the existence and form of solutions.
  An ODE satisfying the condition (continuous coefficients)
  always has solutions, and the solutions form a 2D vector space.
  By Theorem~\ref{thm: linear dependence and Wronskian}
  (\nameref{thm: linear dependence and Wronskian}),
  the linear independence of \(y_1, y_2\) also indicates
  that they form a FSS\@.
\end{remark}

\begin{proof}
  \newcommand{\Sivp}[1]{\mathcal{S}_{\mathtt{IVP}}[#1]}
  \newcommand{\blue}[1]{\textcolor{blue}{#1}}
  \newcommand{\red}[1]{\textcolor{red}{#1}}
  \newcommand{\teal}[1]{\textcolor{teal}{#1}}
  Let \(S\) denote the set of all solutions of the ODE\@.
  We aim to prove that
  \begin{equation*}
    \forall y \in S, \exists! c_1, c_2 \in \mathbb{R} \text{ s.t. }
    y(t) = c_1 y_1(t) + c_2 y_2(t)
    \quad\leftrightarrow\quad
    \exists t_0 \in I \text{ s.t. } W(y_1,y_2)(t_0) \neq 0
  \end{equation*}

  Let \(\Sivp{\teal{t^*},\red{m},\red{n}}\) denote the solution of
  the following IVP
  (by Theorem~\ref{thm: Existence and Uniqueness Theorem},
  there is exactly one solution):
  \begin{equation*}
    y'' + p(t)\cdot y' + q(t)\cdot y = 0
    \qquad y(\teal{t^*}) = \red{m}, \quad y'(\teal{t^*}) = \red{n}
    \qfor t \in I \qand \teal{t^*} \in I
  \end{equation*}

  By Theorem~\ref{thm: Wronskian and IVP with homogeneous ODE},
  we have
  \begin{equation*}
    \forall \teal{t_0} \in I, \forall \red{m},\red{n} \in \mathbb{R}
    \Big[
      \Big(W(y_1,y_2)(\teal{t_0}) \neq 0\Big)
      \; \leftrightarrow \;
      \Big(
        \exists! c_1, c_2 \in \mathbb{R} \text{ s.t. }
        \Sivp{\teal{t_0},\red{m},\red{n}} = c_1 y_1 + c_2 y_2
      \Big)
    \Big]
  \end{equation*}
  Instantiate \(\red{m},\red{n}\) with
  \(\blue{y}(\teal{t_0}), \blue{y'}(\teal{t_0})\)
  for all \(\blue{y} \in S\):
  \begin{equation*}
    \forall \teal{t_0} \in I, \forall \blue{y} \in S \Big[
      \Big(W(y_1,y_2)(\teal{t_0}) \neq 0\Big)
      \; \leftrightarrow \;
      \Big(
        \exists! c_1, c_2 \in \mathbb{R} \text{ s.t. }
        \Sivp{
          \teal{t_0},
          \blue{y}(\teal{t_0}),
          \blue{y'}(\teal{t_0})
        } = c_1 y_1 + c_2 y_2
      \Big)
    \Big]
  \end{equation*}
  Since \(\blue{y}\) satisfies the ODE and initial conditions itself,
  \(\Sivp{\teal{t_0},\blue{y}(\teal{t_0}),\blue{y'}(\teal{t_0})}\)
  is identical to \(\blue{y}\). Thus we obtain
  \begin{equation}
    \forall \teal{t_0} \in I, \forall \blue{y} \in S \Big[
      \Big(W(y_1,y_2)(\teal{t_0}) \neq 0\Big)
      \; \leftrightarrow \;
      \Big(
        \exists! c_1, c_2 \in \mathbb{R} \text{ s.t. }
        \blue{y} = c_1 y_1 + c_2 y_2
      \Big)
    \Big]
    \label{eq: local.03 forall t_0}
  \end{equation}

  Now apply the following two corollaries
  from first-order formal logic:
  \begin{gather*}
    \Big( \exists x. R(x) \Big) \; \land \;
    \forall x. \Big(R(x) \to (P(x)
    \leftrightarrow Q) \Big)
    \quad\implies\quad
    \Big( \exists x. (R(x) \land P(x)) \Big)
    \leftrightarrow Q \\
    \Big( \exists x. R(x) \Big) \; \land \;
    \forall x. \Big(R(x) \to (P(x)
    \leftrightarrow Q) \Big)
    \quad\implies\quad
    \Big( \forall x. (R(x) \to P(x)) \Big)
    \leftrightarrow Q
  \end{gather*}
  Since both \(I\) and \(S\) are non-empty,
  we can transform~\cref{eq: local.03 forall t_0} into
  \begin{equation*}
    \Big( \exists\teal{t_0} \in I,\;
    W(y_1,y_2)(\teal{t_0}) \neq 0\Big)
    \; \leftrightarrow \;
    \forall \blue{y} \in S \Big(
      \exists! c_1, c_2 \in \mathbb{R} \text{ s.t. }
      \blue{y} = c_1 y_1 + c_2 y_2
    \Big)
  \end{equation*}
  This completes the proof.
\end{proof}

\begin{theorem}[Existence of Fundamental Set of Solutions]
  {}\label{thm: Existence of Fundamental Set of Solutions}
  If \(p(t)\) and \(q(t)\) are continuous on an open interval \(I\),
  then there must exist fundamental sets of solutions for the following ODE\@:
  \begin{equation}
    y'' + p(t)\cdot y' + q(t)\cdot y = 0
    \qfor t \in I
    \label{eq: local.04 ODE}
  \end{equation}

  A FSS can be constructed as follows:
  Pick any \(t_0 \in I\), and let \(y_1\) be the solution
  of \cref{eq: local.04 ODE} with initial conditions
  \begin{equation*}
    y_1(t_0) = 1, \quad y_1'(t_0) = 0
  \end{equation*}
  Let \(y_2\) be the solution of \cref{eq: local.04 ODE} with initial conditions
  \begin{equation*}
    y_2(t_0) = 0, \quad y_2'(t_0) = 1
  \end{equation*}
  Then, \(y_1\) and \(y_2\) form a FSS of \cref{eq: local.04 ODE}.
\end{theorem}

\begin{proof}
  Theorem~\ref{thm: Existence and Uniqueness Theorem}
  (\nameref{thm: Existence and Uniqueness Theorem})
  guarantees the existence of \(y_1\) and \(y_2\).
  By the definition of Wronskian, we have
  \begin{equation*}
    W(y_1,y_2)(t_0) =
    \begin{vmatrix}
      1 & 0 \\
      0 & 1
    \end{vmatrix} = 1 \neq 0
  \end{equation*}
  By the \hyperref[thm: corollary of Abel's Theorem]{
  corollary of Abel's Theorem},
  \(W(y_1,y_2)(t) \neq 0\) for all \(t \in I\).
  \(y_1\) and \(y_2\) form a FSS of \cref{eq: local.04 ODE}.
\end{proof}

\section{Homogeneous Equations with Constant Coefficients}

\subsection{}

\begin{theorem}
  {}\label{thm: complex solution to real solutions}
  Given an ODE with \textbf{real-valued} coefficients
  \begin{equation*}
    y'' + p(t)\cdot y' + q(t)\cdot y = 0
  \end{equation*}
  if \(y\) is a \textbf{complex-valued} solution of the ODE\@,
  then \(\Re(y), \Im(y)\) are also solutions of the ODE\@.
\end{theorem}

\begin{proof}
  Substitute \(y(t) = u(t) + \iu \, v(t)\) into the ODE\@,
  and it can be verified that both \(u(t)\) and \(v(t)\)
  satisfy the ODE\@.
\end{proof}

\begin{theorem}
  For the linear ODE with constant coefficients
  \begin{equation}
    a y'' + b y' + c y = 0
  \end{equation}
  the \textbf{characteristic equation} is given by
  \begin{equation*}
    a r^2 + b r + c = 0
  \end{equation*}
  Let \(r_1, r_2\) be the two roots and \(\Delta := b^2 - 4 a c\)
  be the discriminant. Then,
  \begin{itemize}
    \item If \(\Delta > 0\), then \(r_1 \neq r_2 \in \mathbb{R}\),
      and the general solution is given by
      \begin{equation*}
        y(t) = c_1 e^{r_1 t} + c_2 e^{r_2 t}
      \end{equation*}
    \item If \(\Delta = 0\), then \(r_1 = r_2 \in \mathbb{R}\),
      and the general solution is given by
      \begin{equation*}
        y(t) = (c_1 + c_2 t) e^{-\frac{b}{2 a} t}
      \end{equation*}
    \item If \(\Delta < 0\), then \(r_1 = \overline{r_2} \in \mathbb{C}\).
      Let \(\lambda := \Re(r_1), \quad \mu := \Im(r_1)\).
      The general solution is given by
      \begin{gather*}
        y(t) = e^{\lambda t} \Big(
          c_1 \cos(\mu t) + c_2 \sin(\mu t)
        \Big)
        \qquad \text{(real-valued when \(a,b,c,c_1,c_2\in\mathbb{R}\))} \\
        \qor
        y(t) = e^{\lambda t} \Big(
          c_1 \cos(\mu t) + \textcolor{red}{\iu}\, c_2 \sin(\mu t)
        \Big)
        \qquad \text{(complex-valued)}
      \end{gather*}
  \end{itemize}
\end{theorem}

\begin{proof}
  The general solutions are derived by substituting
  the ansatz \( y(t) = e^{rt} \) into the differential equation
  and applying
  Theorem~\ref{thm: Wronskian and general solution of a homogeneous ODE}.

  When \(\Delta < 0\),
  the characteristic equation yields complex conjugate roots,
  leading to complex-valued solutions.
  However, if the coefficients \( a, b, c \in \mathbb{R} \),
  we may apply Theorem~\ref{thm: complex solution to real solutions}
  to construct an equivalent real-valued general solution.

  When \(\Delta = 0\),
  he characteristic equation has a repeated real root \( r \),
  yielding only one solution \( y_1(t) = e^{rt} \).
  To form a fundamental set of solutions,
  we require a second linearly independent solution.
  Using Theorem~\ref{thm: Abel's Theorem} (\nameref{thm: Abel's Theorem}),
  we can obtain \( y_2(t) = t e^{rt} \).
  \textit{(Note: Any solution linearly independent of \( y_1(t) \) is
  acceptable, though \( t e^{rt} \) is the standard choice.)}
\end{proof}

\begin{theorem}[Reduction of Order]
  {}\label{thm: Reduction of Order}
  If \(y_1\) is a non-zero solution of the ODE
  \begin{equation*}
    y'' + p(t)\cdot y' + q(t)\cdot y = 0
    \qfor t \in I
  \end{equation*}
  then a second linearly independent solution \(y_2\) can be found by
  \begin{equation*}
    y_2(t) = y_1(t) \int
    \frac{\exp(-\int p(t) \dd{t})}{{[y_1(t)]}^2} \dd{t}
  \end{equation*}
\end{theorem}

\begin{proof}
  We suppose \(y_2(t) = v(t) y_1(t)\) for some function \(v\).
  Substituting \(y_2\) into the ODE\@, we have
  \begin{equation*}
    y_1 v'' + (2 y_1' + p y_1) v' = 0
  \end{equation*}
  Thus \(v\) can be solved.
  \begin{equation*}
    v'(t) = \frac{1}{y_1^2(t)} \exp(-\int p(t) \dd{t})
  \end{equation*}
\end{proof}

\subsection{Euler Equations}

\begin{definition}[Euler Equation]
  {}\label{def: Euler Equation}
  Euler equations are the differential equations of the form
  \begin{equation}
    x^2 \dv[2]{y}{x} + A x \dv{y}{x} + B y = 0
    \qfor x > 0
    \label{eq: Euler Equation}
  \end{equation}
  where \(A, B\) are constants.
\end{definition}

Let \(x = e^t\), then we have
\begin{align*}
  \dv{y}{x} &= \dv{y}{(e^t)}
  &
  \dv[2]{y}{x}&= \dv{(e^t)} \left( \dv{y}{x} \right) \\
  &= \frac{1}{e^t} \cdot \dv{y}{t}
  &
  &= \frac{1}{e^t \dd{t}} \cdot \dd( \frac{1}{e^t} \dv{y}{t} ) \\
  &&&= \frac{1}{e^t \dd{t}} \left[
    - \frac{\dd{t}}{e^t} \dv{y}{t} + \frac{1}{e^t} \dd(\dv{y}{t})
  \right] \\
  &&&= \frac{1}{e^{2 t}} \left[
    - \dv{y}{t} + \dv[2]{y}{t}
  \right]
  \intertext{Thus,}
  x \dv{y}{x} &= \dv{y}{t}
  &
  x^2 \dv[2]{y}{x} &= \dv[2]{y}{t} - \dv{y}{t}
\end{align*}
Substituting the above into \cref{eq: Euler Equation}, we have
\begin{equation*}
  \dv[2]{y}{t} + (A - 1) \dv{y}{t} + B y = 0
\end{equation*}
which is a linear homogeneous ODE with constant coefficients.

\section{Non-Homogeneous Equations}

\begin{theorem}
  {}\label{thm: general solution of non-homogeneous ODE}
  Suppose \(p(t), q(t), r(t)\) are continuous on an interval \(I\).
  Then the \textbf{general solution} of the ODE
  \begin{equation*}
    y'' + p(t)\cdot y' + q(t)\cdot y = r(t)
  \end{equation*}
  is given by
  \begin{equation*}
    y(t) = y_h(t) + y_p(t)
  \end{equation*}
  where \(y_h\), or denoted as \(y_c\) and called
  the \textbf{complementary solution},
  is the general solution of the corresponding homogeneous ODE
  \begin{equation*}
    y'' + p(t)\cdot y' + q(t)\cdot y = 0
  \end{equation*}
  and \(y_p\) is a \textbf{particular solution} of the non-homogeneous ODE\@.
\end{theorem}


\subsection{Method of Undetermined Coefficients}

Consider the non-homogeneous ODE with constant coefficients
\begin{equation*}
  a y'' + b y' + c y = r(t)
\end{equation*}
When the non-homogeneous term \(r(t)\) is of a specific form,
we can apply the method of undetermined coefficients to find
a particular solution \(y_p\).
And since \(a y'' + b y' + c y = 0\) can be easily solved,
we can obtain the general solution of the non-homogeneous ODE\@.

\begin{center}
  \newcommand{\red}[1]{\textcolor{red}{#1}}
  \newcommand{\blue}[1]{\textcolor{blue}{#1}}
  \newcommand{\teal}[1]{\textcolor{teal}{#1}}
  \newcommand{\brown}[1]{\textcolor{brown}{#1}}
  \begin{tabular}{c @{\hspace{1cm}} c}
    \toprule
    \(r(t)\) & \(y_p(t)\) \\
    \midrule
    \(a_n t^n + \cdots + a_{1} t + a_0\) &
    \(\red{t^s} (A_n t^n + \cdots + A_{1} t + A_0)\) \\[1em]
    \((a_n t^n + \cdots + a_{1} t + a_0) \blue{e^{\alpha t}}\) &
    \(\red{t^s} (A_n t^n + \cdots + A_{1} t + A_0) \blue{e^{\alpha t}}\) \\[1em]
    \(
      \begin{matrix}
        (a_n t^n + \cdots + a_{1} t + a_0) \blue{e^{\alpha t}}
        \teal{\sin\beta t} \\
        (a_n t^n + \cdots + a_{1} t + a_0) \blue{e^{\alpha t}}
        \teal{\cos\beta t}
    \end{matrix} \text{ or}\) &
    \(\red{t^s} \Big[
        (A_n t^n + \cdots + A_{1} t + A_0) \teal{\sin\beta t} +
        (B_n t^n + \cdots + B_{1} t + B_0) \teal{\cos\beta t}
    \Big] \blue{e^{\alpha t}} \) \\[1em]
    \bottomrule
  \end{tabular}

  \smallskip
  \begin{minipage}{0.9\textwidth}
    \textit{Notes.}
    Here, \(\red{s}\) is the smallest non-negative integer
    (\(\red{s} \in \{0, 1, 2\} \)) ensuring that
    \(y_p(t)\) and \(y_h(t)\) possess no common terms.
    Equivalently, \(\red{s}\) represents the multiplicity of the root
    \(\brown{0}\), \(\brown{\alpha}\), or \(\brown{\alpha + i\beta}\)
    in the characteristic equation,
    corresponding to the respective forms of
    the non-homogeneous term \(r(t)\).
  \end{minipage}
\end{center}

\begin{theorem}
  If \(y_1\) is a solution to \(a y'' + b y' + c y = r_1(t)\)
  and \(y_2\) is a solution to \(a y'' + b y' + c y = r_2(t)\),
  then \(y_1 + y_2\) is a solution to the equation
  \(a y'' + b y' + c y = r_1(t) + r_2(t)\).
\end{theorem}

\subsection{Variation of Parameters}

\begin{theorem}[Variation of Parameters]
  {}\label{thm: Variation of Parameters}
  Suppose \(p(t), q(t), r(t)\) are continuous on an open interval \(I\).
  Then a particular solution of the ODE
  \begin{equation}
    y'' + p(t)\cdot y' + q(t)\cdot y = r(t)
    \qfor t \in I
    \label{eq: local.07 non-homogeneous CC ODE}
  \end{equation}
  can be found by
  \begin{equation}
    y_p(t) = u_1(t) y_1(t) + u_2(t) y_2(t)
    \label{eq: local.08 a particular solution}
  \end{equation}
  where \((y_1, y_2)\) is a fundamental set of solutions of the
  corresponding homogeneous ODE \(y'' + p(t)\cdot y' + q(t)\cdot y = 0\),
  and \(u_1, u_2\) are functions satisfying
  \begin{equation}
    \begin{pmatrix}
      y_1 & y_2 \\ y_1' & y_2'
    \end{pmatrix}
    \begin{pmatrix}
      u_1' \\ u_2'
    \end{pmatrix} =
    \begin{pmatrix}
      0 \\ r
    \end{pmatrix}
    \label{eq: local.09 conditions for u_1 and u_2}
  \end{equation}
  And the general solution is given by
  \begin{equation}
    y(t) = (u_1(t) + c_1) y_1(t) + (u_2(t) + c_2) y_2(t)
    \label{eq: local.11 the general solution}
  \end{equation}
  for constants \(c_1, c_2\).
\end{theorem}

In fact, \(u_1, u_2\) can be given by
\begin{equation}
  u_1(t) = - \int \frac{y_2(t) r(t)}{W(y_1,y_2)(t)} \dd{t}
  \qquad\qand\qquad
  u_2(t) = \int \frac{y_1(t) r(t)}{W(y_1,y_2)(t)} \dd{t}
  \label{eq: local.10 result of u_1 and u_2}
\end{equation}

\begin{proof}[Variation of Parameters]
  Given \cref{eq: local.08 a particular solution} and
  \cref{eq: local.09 conditions for u_1 and u_2},
  we can verify that \(y_p\) is a particular solution of
  the non-homogeneous ODE~\ref{eq: local.07 non-homogeneous CC ODE}\@.
  Since \(y_1, y_2\) form a FSS of the corresponding homogeneous ODE,
  \(W(y_1, y_2)(t) \neq 0\) for all \(t \in I\)
  (by Theorem~\ref{thm: Wronskian and general solution of a homogeneous ODE}
  and the \hyperref[thm: Abel's Theorem]{corollary of Abel's Theorem}).
  Thus, \(u_1', u_2'\) can be solved by Cramer's rule from
  \cref{eq: local.09 conditions for u_1 and u_2}.
  After integration, \(u_1, u_2\) in \cref{eq: local.10 result of u_1 and u_2}
  is obtained and so is the particular solution \(y_p\) in
  \cref{eq: local.08 a particular solution}.
  Finally, using Theorem~\ref{thm: general solution of non-homogeneous ODE},
  we have the general solution in \cref{eq: local.11 the general solution}.
\end{proof}



\end{document}
